\section{Local representatives of the equations of motion}
\label{ch02:equation-section}

The ordinary section $P_d^{\rm loc}$ reproduces Taylor coefficients,
but need not reproduce equations of motion. We construct a second
section for this purpose. Everything in this section concerns the
finite-volume Wilson measure; no renormalization estimate is used.

A radial polynomial potential $A_j$, homogeneous of degree $j+1$,
is a \emph{Maxwell solution potential} if
\[
 \sum_\mu\partial_\mu(dA_j)_{\mu\nu}=0\quad\hbox{for every }\nu.
\]
Its curvature is automatically closed. Let $\mathcal I_d\subset V_d$
be the polynomials vanishing on all such solution jets through degree
$d$. This is the degree-$d$ part of the equation ideal; multiplication
is understood only when the resulting scaling dimension remains at most $d$.
The condition preserves both scaling dimension and factor count.
It is the kernel of restriction to solution jets, rather than a
statement about the interacting measure.

\begin{theorem}[An equation-adapted local section]
\label{ch02:equation-adapted}
There are a real, symmetry-equivariant linear section
$\mathcal V_d$ from $V_d$ to local continuous gauge-invariant functions of links, a real equivariant automorphism
$R_d$ of $V_d$, and finite constants $r_d,M_d$, with the following
properties:
\begin{enumerate}
\item $\mathcal V_d1=1$, $R_d1=1$, and
$\operatorname{Tay}_d\mathcal V_d p=p$.
The section has a fixed finite stencil, independent of the torus.
\item If $p$ has factor count $s$ and scaling dimension $w$,
then $R_dp-p$ is a sum of terms whose pair
$(\text{factor count},\text{scaling dimension})$ is strictly larger than $(s,w)$
in lexicographic order. Every term of $\mathcal V_dp$ is a product
of at least $s$ transported curvature coordinates.
\item For every $p\in\mathcal I_d$, every torus of side at least $M_d$,
every anchor $o$, every real $\beta\ne0$, and every bounded measurable
$S$ depending only on links with both endpoints farther than $r_d$
from $o$,
\begin{equation}
 \int \mathcal V_d(R_dp)(o,U)S(U)\,d\mu_\beta(U)=0.
 \label{ch02:microscopic-null}
\end{equation}
No gauge invariance or regularity of $S$ is assumed.
\end{enumerate}
The small-field derivative bounds for a fixed section proved above
also apply to $\mathcal V_d$, with its own finite constants.
\end{theorem}

The regrading $R_d$ is essential: the exact interacting relation is
$\mathcal V_d(R_d\mathcal I_d)$, not an assertion that
$P_d^{\rm loc}(\mathcal I_d)$ has zero Wilson correlations.
The proof below supplies the finite-stencil construction, whereas the
compact-group integration by parts underlying it is the usual lattice
equation-of-motion identity for Wilson's measure
\cite{Wilson1974,OsterwalderSeiler1978}.

\subsection{Curvature tests avoiding a prescribed link}

For $u\in\gLie$, pairing the transported curvature
\eqref{ch02:transported-curvature} with $u$ gives the real-valued function
$U\mapsto\langle u,\sigma_{\gamma,\square}(U)\rangle_{\gLie}$.
We call any such function a curvature coordinate.
Call it admissible for an oriented link $e$ if neither its plaquette
nor its transporting path uses the underlying unoriented link $e$.
Such a coordinate is independent of $U_e$.

\begin{lemma}\label{ch02:avoiding-tests}
For each $e$, the linearizations at the zero jet of admissible
curvature coordinates span the dual of the finite radial jet space
$\mathcal J_{\le d}$.
\end{lemma}
\begin{proof}
It suffices to show that their common annihilator is zero.
Removing one link does not disconnect $\ZZ^4$, so every plaquette
not containing $e$ admits a transporting path avoiding $e$.
Let $A$ be the polynomial one-form representing a jet in the annihilator.
At the zero connection the derivative of transport around an oriented
unit plaquette $\square$ in direction $A$ is
$\oint_{\partial\square}A=\int_\square dA$ by Stokes' theorem.
Transport to the anchor does not change this first variation.
Consequently $\int_\square dA=0$ for every plaquette that does not
contain $e$.

For each orientation $(\mu,\nu)$, this integral is a polynomial in
the base point $z$ of the plaquette. It vanishes at all but finitely
many integer points and hence identically, by the one-variable polynomial
identity applied successively in each coordinate. Let $A_j$ be the
highest nonzero homogeneous layer of $A$. Averaging a polynomial over
a translated unit square preserves its highest homogeneous part as a
polynomial in $z$. Thus the degree-$j$ part of the integral is
$(dA_j)_{\mu\nu}(z)$. Since the integral vanishes for every orientation,
$dA_j=0$. Radial gauge and Euler's identity give
\[
 0=\partial_i(x^\mu A_{j,\mu})
   =A_{j,i}+x^\mu\partial_i A_{j,\mu}
   =(j+2)A_{j,i},
\]
where closedness was used in the last equality.
This contradicts the choice of $j$. Descending through the layers
proves the assertion. Finite-dimensional duality now selects a
finite spanning collection of admissible coordinates.
\end{proof}

\subsection{Construction of the equation-adapted section}

\begin{proof}[Proof of Theorem~\ref{ch02:equation-adapted}]
We first represent the generators of the equation ideal by Wilson forces.
We then integrate by parts in a single link and assemble the local section.

\smallskip\noindent\emph{Step 1. Representation by Wilson forces.}

Write $S_W$ for the Wilson action without its factor $\beta$.
For a positive link $e$ and a direction $X\in\gLie$,
the force $F_e\in\gLie$ is characterized by
\[
 \partial_{e,X}S_W=\langle X,F_e\rangle,\qquad
 \partial_{e,X}f(U)=
 \left.\frac{d}{dt}f(U_e\leftarrow e^{tX}U_e)\right|_{t=0}.
\]
Differentiating the plaquettes incident on $e$ expresses $F_e$
as a signed sum of plaquette-curvature coordinates based at the
initial vertex of $e$. The precise signs follow the chosen
plaquette orientations. Transport the force to the anchor by a
path avoiding $e$, and pair it there with an arbitrary element of $\gLie$.
At the zero connection the force is the signed sum, to first order in
$A$, of the integrals $\int_\square dA$ over plaquettes incident on $e$.

Fix a layer $j$. These linear force functionals, restricted to jets
through layer $j$, span every functional which acts only on $A_j$
and annihilates the Maxwell solution potentials in that layer.
Indeed, if all force functionals vanish, their values at integer
base points form identically zero polynomials. Their top-degree
terms are
\[
 -\sum_\mu\partial_\mu(dA_j)_{\mu\nu}.
\]
For a link in direction $\nu$, the two incident $(\mu,\nu)$ plaquettes
enter with opposite signs. Their integrals therefore give a first
difference in the $\mu$ direction. Thus the top layer is a Maxwell
solution potential. Every functional of the indicated type
annihilates the common force kernel and is therefore a finite
linear combination of these forces. Notice that this argument
identifies the top layer only; it does not identify the entire
kernel of the discrete force system with the continuum solution
jet space.

Here is the corresponding description of the polynomial ideal.
Choose a linear projection of each jet layer onto its solution
subspace and let $\mathcal E$ be their direct sum. In the full,
not necessarily invariant, polynomial algebra the kernel of
$p\mapsto p(\mathcal EJ)$ is generated by linear functionals
vanishing on those solution subspaces. To check this without
an ideal-generation theorem, use
\[
 \prod_{\nu=1}^s u_\nu-\prod_{\nu=1}^s(\mathcal E^*u_\nu)
 =\sum_{i=1}^s
 (u_i-\mathcal E^*u_i)
 \prod_{\nu<i}\mathcal E^*u_\nu
 \prod_{\nu>i}u_\nu .
\]
Choose homogeneous linear coordinates, so both gradings are
preserved. Every homogeneous element of $\mathcal I_d$ is
therefore a sum of products $\lambda z_1\cdots z_{s-1}$,
where $\lambda$ is a one-layer equation functional.
Individual summands need not be gauge invariant.

Represent $\lambda$ by the finite force combination just constructed
on its truncated jet space. For each force summand involving $e$,
represent each $z_i$ by a linear combination of curvature coordinates
avoiding $e$, using Lemma~\ref{ch02:avoiding-tests} on the \emph{full}
jet space through degree $d$. Multiply these representatives by
that force summand. The resulting function $f$ has at least $s$
curvature factors. Its Taylor polynomial is
\begin{equation}
 \operatorname{Tay}_d f
   =\lambda z_1\cdots z_{s-1}
       +\{\hbox{terms strictly higher in }(s,w)\}.
 \label{ch02:equation-triangle}
\end{equation}
There are two distinct errors. The nonlinear part of a curvature
or force coordinate increases factor count. The linear force
representation was fixed only through the layer of $\lambda$;
its remaining linear terms increase the scaling dimension while preserving the factor
count. This is why the order in
\eqref{ch02:equation-triangle} puts factor count first.

\smallskip\noindent\emph{Step 2. Integration by parts and construction of the section.}

Every force term above has the form
$H\langle X,F_e\rangle$, where $X$ is the fixed direction at the
anchor transported back to the link, and $H$ is the product of the
remaining curvature coordinates and the exterior spectator.
Freeze the other links. Both $X$ and $H$ are then constant in $U_e$. Haar invariance
on the compact group gives
\[
 0=\int_G\partial_{e,X}\{e^{-\beta S_W}\}\,dU_e
   =-\beta\int_G\langle X,F_e\rangle
                          e^{-\beta S_W}\,dU_e.
\]
Multiplication by $H$ is legitimate because it is constant in
$U_e$. Fubini then proves the desired zero expectation, even for a
merely bounded measurable exterior spectator. This proof never
differentiates that spectator.

All lists used in constructing a basis are finite. Enlarge one
radius to contain their plaquettes, paths, and differentiated links.
Enlarge one period threshold so that distinct relevant cover links
cannot become identified with a differentiated link after
periodization. These two choices work for every anchor by translation.
They are independent of $\beta$ and of the exterior spectator.

Average the constructed functions over the gauge element at the
anchor. Every transported curvature transforms there by the adjoint
action, so this is a finite-dimensional average of tensors under the adjoint action.
It preserves the number of curvature factors and
\eqref{ch02:equation-triangle}. It also preserves the exact zero
identity: gauge invariance of the Wilson law transfers the gauge
transformation to the bounded exterior spectator without changing
its support. The result is fully gauge invariant.

Choose a basis of each homogeneous part of $\mathcal I_d$, perform
this construction there, and extend to a basis of $V_d$ by using
$P_d^{\rm loc}$ on a complement. The resulting linear map
$\mathcal F_d$ preserves the unit. Its Taylor map is
\[
 R_d=\operatorname{Tay}_d\mathcal F_d=I+T_d,
\]
where $T_d$ strictly increases the finite lexicographic grading.
Thus $T_d$ is nilpotent and
$R_d^{-1}=I-T_d+T_d^2-\cdots$ is a finite sum.
Average $\mathcal F_d$ simultaneously on domain and range over
signed coordinate permutations and over conjugation. These
operations preserve the ideal, both gradings, and the unit;
hence the same triangular conclusion remains valid.
Finally put
\[
 \mathcal V_d=\mathcal F_dR_d^{-1}.
\]
This gives the Taylor identity, all symmetries, and
$\mathcal V_dR_d p=\mathcal F_dp$ for $p\in\mathcal I_d$.
Both $R_d$ and its inverse preserve the lower bound on factor
count. The fixed finite products therefore obey exactly the
small-field derivative argument used for $P_d^{\rm loc}$,
with new constants for the new stencil.
\end{proof}

The identity \eqref{ch02:microscopic-null} is exact at each cutoff.
It gives no bound on a renormalized coefficient change.
To use it for continuum null fields one must compare this section
with the section used along the actual renormalization history,
with one invertible comparison map for all labels and all spectators.
